\section{Monte Carlo Design}

\subsection{Direct-inversion benchmark}

The first experiment fixes the baseline model before introducing WGAN-GP or
DML. Values are lognormal with $\log V_{ij}\sim\mathcal N(0,0.5^2)$, every
auction has $N=5$ bidders, and the econometrician observes only the maximum.
For each sample size $n\in\{200,500,2000,10000\}$, I generate 500 Monte Carlo
samples. The random seed is fixed at 20260716. In every replication,
$\widehat F_{\mathrm{DI}}=\widehat H_n^{1/N}$ and the reserve maximizes
$r\{1-\widehat F_{\mathrm{DI}}(r)\}$ over the observed maxima.

\begin{table}[htbp]
\centering
\caption{Direct-inversion reserve estimates}
\label{tab:baseline-direct-inversion}
\begin{tabular}{rrrrrr}
\toprule
$n$ & Mean $\widehat r$ & Bias & SD & RMSE & Mean regret \\
\midrule
200 & 0.9053 & 0.1335 & 0.0794 & 0.1552 & 0.0032 \\
500 & 0.8446 & 0.0728 & 0.0635 & 0.0965 & 0.0009 \\
2000 & 0.7960 & 0.0242 & 0.0556 & 0.0606 & 0.0002 \\
10000 & 0.7584 & -0.0134 & 0.0513 & 0.0530 & 0.0001 \\
\bottomrule
\end{tabular}

\begin{minipage}{0.94\textwidth}
\footnotesize
\emph{Notes:} The true reserve solves
$1-F_0(r)-rf_0(r)=0$. Revenue regret is
$\mathcal R_N(r_0;F_0)-\mathcal R_N(\widehat r;F_0)$.
\end{minipage}
\end{table}

\begin{figure}[htbp]
\centering
\includegraphics[width=\textwidth]{baseline_direct_inversion.png}
\caption{Recovery of the latent CDF and sampling distribution of the
direct-inversion reserve estimator}
\label{fig:baseline-direct-inversion}
\end{figure}

This experiment supplies a numerical identification check and a calibration
target for every later learner. The left panel of
Figure~\ref{fig:baseline-direct-inversion} shows where the inversion has little
information; the right panel and Table~\ref{tab:baseline-direct-inversion}
measure whether this translates into economically material reserve error.

\subsection{Planned estimator comparison}

The next experiments will vary the auction sample size, number of bidders,
valuation distribution, degree of tail thickness, auction-level heterogeneity,
and structural misspecification. Candidate valuation distributions include
lognormal, gamma, mixture, and nonregular designs. The comparison will include
direct order-statistic inversion, conventional plug-in estimators, WGAN-GP
plug-in estimation, cross-fitted WGAN-GP, and the orthogonal WGAN-DML
estimator. Evaluation metrics include bias, root mean squared error,
confidence-interval coverage, distributional fit, and revenue regret:
\begin{equation}
R(r_0;F_0)-R(\widehat r;F_0).
\end{equation}
The same Monte Carlo seed policy and the direct benchmark will be retained
when these estimators are added.
